% Supplemental mathematical appendix. Include only if the venue permits it;
% the main paper remains the canonical submission body.
\section{Algebra and randomisation geometry of PLACEBO}
\label{app:placebo-math}

\paragraph{Contrast completeness.}
For the contrast matrix $C$ in Section~\ref{sec:inference}, direct row reduction
gives $\operatorname{rank}(C)=4$ and $C\mathbf 1_5=0$. Hence
$\ker(C)=\operatorname{span}\{\mathbf 1_5\}$ and the registered rows span every
identified contrast among the five potential-outcome means. The two displayed
decomposition identities follow by adding the corresponding rows of $C$.

\paragraph{Orthogonal coordinates.}
The coefficient vectors of $(s,p,g,d)$ are pairwise orthogonal and each sums to
zero. They are nonzero, so they form an orthogonal basis of the four-dimensional
contrast space $\mathbf 1_5^\perp$. This basis isolates specificity, the
packet-bearing bundle, continuation gain, and duplicate-control disagreement
without asserting a factorial decomposition that the four arms cannot identify.

\paragraph{Allocation orbit.}
Before outcome observation, the four labelled slots carry a uniformly random
permutation of $(R,S,N,Z)$. Because $N$ and $Z$ receive the same intervention,
the scientific allocation has $4!/2!=12$ elements: the cosets of the subgroup
$\langle(N\ Z)\rangle$ in $S_4$. This is a coset space, not a quotient group,
because that subgroup is not normal. The independent fair label swap restores
the full labelled assignment needed for the duplicate-control audit.

\paragraph{Invariant and anti-invariant components.}
Let $T$ swap $N$ and $Z$. Then $T$ fixes $s,p,g$ and maps $d$ to $-d$.
Conditionally on the observed magnitudes $|D_i|$, exchangeability makes the
audit statistic $\sum_i w_i\epsilon_iD_i$ an exact Rademacher randomisation
variable. For equal weights $w_i=1/N$ and binary outcomes, exactly $m$
discordant pairs give
\[
q_0=\frac{m}{N},\qquad
\operatorname{Var}_{\epsilon}\!\left(\frac1N\sum_i\epsilon_iD_i
\middle| |D|\right)=\frac{m}{N^2}=\frac{q_0}{N}.
\]

\paragraph{Intersection--union decision.}
The scientific alternative is $H_1=\{\tau_{\rm content}>0\}\cap
\{\tau_{\rm excess}>0\}$. Its null is the union of the two component nulls.
Rejecting only when both level-$\alpha$ component tests reject has size at most
$\alpha$: under any point in the union, at least one component null is true,
and joint rejection is a subset of that true-null test's rejection event.
